% !TeX root = ../../Book5.tex
% 纲要标题：### 23.4 任意基域下的曲线对偶
% 纲要位置：几何学计划纲要.md，第 5972 行。
% 正文状态：未撰写；本文件目前是框架。

% 纲要内容开始
% 
% 本节保持任意基域的 Riemann–Roch 目标；对偶工具取本卷已经建立的代数方法，不引用复解析曲线结论。
% 
% * 光滑曲线的典范线丛与迹
% 
%   把第22章的代数构造专门化到曲线，重新写出配对和迹映射，明确其独立于所选射影嵌入。
% 
% * 曲线对偶的非退化性
% 
%   在曲线的有限消解和长正合列中核验非退化性及自然性，保留非代数闭基域的线性代数。
% 
% * 不可分闭点与坐标留数的限制
% 
%   几何光滑不意味着所有闭点的剩余域扩张都可分；不把分离变量后的系数域公式无条件用于这些点。
% 
% * 点处修改与闭点次数
% 
%   以商层的 k 维数得到 [κ(x):k] 的贡献，为下一节的统一次数公式准备。
% 
% 纲要内容结束

\section{任意基域下的曲线对偶}
\label{g5:sec:2304}

待撰写

\subsection{光滑曲线的典范线丛与迹}
\label{g5:subsec:2304:01}
待撰写

\subsection{曲线对偶的非退化性}
\label{g5:subsec:2304:02}
待撰写

\subsection{不可分闭点与坐标留数的限制}
\label{g5:subsec:2304:03}
待撰写

\subsection{点处修改与闭点次数}
\label{g5:subsec:2304:04}
待撰写
% 纲要条目排为编号小节标题，并进入目录与书签；说明段仍保留于纲要注释中。
